% Options for packages loaded elsewhere
\PassOptionsToPackage{unicode}{hyperref}
\PassOptionsToPackage{hyphens}{url}
\documentclass[
  polish,
  11pt,
  a4paper,
]{article}
\usepackage{xcolor}
\usepackage[margin=2cm]{geometry}
\usepackage{amsmath,amssymb}
\setcounter{secnumdepth}{-\maxdimen} % remove section numbering
\usepackage{iftex}
\ifPDFTeX
  \usepackage[T1]{fontenc}
  \usepackage[utf8]{inputenc}
  \usepackage{textcomp} % provide euro and other symbols
\else % if luatex or xetex
  \usepackage{unicode-math} % this also loads fontspec
  \defaultfontfeatures{Scale=MatchLowercase}
  \defaultfontfeatures[\rmfamily]{Ligatures=TeX,Scale=1}
\fi
\usepackage{lmodern}
\ifPDFTeX\else
  % xetex/luatex font selection
\fi
% Use upquote if available, for straight quotes in verbatim environments
\IfFileExists{upquote.sty}{\usepackage{upquote}}{}
\IfFileExists{microtype.sty}{% use microtype if available
  \usepackage[]{microtype}
  \UseMicrotypeSet[protrusion]{basicmath} % disable protrusion for tt fonts
}{}
\makeatletter
\@ifundefined{KOMAClassName}{% if non-KOMA class
  \IfFileExists{parskip.sty}{%
    \usepackage{parskip}
  }{% else
    \setlength{\parindent}{0pt}
    \setlength{\parskip}{6pt plus 2pt minus 1pt}}
}{% if KOMA class
  \KOMAoptions{parskip=half}}
\makeatother
\usepackage{color}
\usepackage{fancyvrb}
\newcommand{\VerbBar}{|}
\newcommand{\VERB}{\Verb[commandchars=\\\{\}]}
\DefineVerbatimEnvironment{Highlighting}{Verbatim}{commandchars=\\\{\}}
% Add ',fontsize=\small' for more characters per line
\usepackage{framed}
\definecolor{shadecolor}{RGB}{248,248,248}
\newenvironment{Shaded}{\begin{snugshade}}{\end{snugshade}}
\newcommand{\AlertTok}[1]{\textcolor[rgb]{0.94,0.16,0.16}{#1}}
\newcommand{\AnnotationTok}[1]{\textcolor[rgb]{0.56,0.35,0.01}{\textbf{\textit{#1}}}}
\newcommand{\AttributeTok}[1]{\textcolor[rgb]{0.13,0.29,0.53}{#1}}
\newcommand{\BaseNTok}[1]{\textcolor[rgb]{0.00,0.00,0.81}{#1}}
\newcommand{\BuiltInTok}[1]{#1}
\newcommand{\CharTok}[1]{\textcolor[rgb]{0.31,0.60,0.02}{#1}}
\newcommand{\CommentTok}[1]{\textcolor[rgb]{0.56,0.35,0.01}{\textit{#1}}}
\newcommand{\CommentVarTok}[1]{\textcolor[rgb]{0.56,0.35,0.01}{\textbf{\textit{#1}}}}
\newcommand{\ConstantTok}[1]{\textcolor[rgb]{0.56,0.35,0.01}{#1}}
\newcommand{\ControlFlowTok}[1]{\textcolor[rgb]{0.13,0.29,0.53}{\textbf{#1}}}
\newcommand{\DataTypeTok}[1]{\textcolor[rgb]{0.13,0.29,0.53}{#1}}
\newcommand{\DecValTok}[1]{\textcolor[rgb]{0.00,0.00,0.81}{#1}}
\newcommand{\DocumentationTok}[1]{\textcolor[rgb]{0.56,0.35,0.01}{\textbf{\textit{#1}}}}
\newcommand{\ErrorTok}[1]{\textcolor[rgb]{0.64,0.00,0.00}{\textbf{#1}}}
\newcommand{\ExtensionTok}[1]{#1}
\newcommand{\FloatTok}[1]{\textcolor[rgb]{0.00,0.00,0.81}{#1}}
\newcommand{\FunctionTok}[1]{\textcolor[rgb]{0.13,0.29,0.53}{\textbf{#1}}}
\newcommand{\ImportTok}[1]{#1}
\newcommand{\InformationTok}[1]{\textcolor[rgb]{0.56,0.35,0.01}{\textbf{\textit{#1}}}}
\newcommand{\KeywordTok}[1]{\textcolor[rgb]{0.13,0.29,0.53}{\textbf{#1}}}
\newcommand{\NormalTok}[1]{#1}
\newcommand{\OperatorTok}[1]{\textcolor[rgb]{0.81,0.36,0.00}{\textbf{#1}}}
\newcommand{\OtherTok}[1]{\textcolor[rgb]{0.56,0.35,0.01}{#1}}
\newcommand{\PreprocessorTok}[1]{\textcolor[rgb]{0.56,0.35,0.01}{\textit{#1}}}
\newcommand{\RegionMarkerTok}[1]{#1}
\newcommand{\SpecialCharTok}[1]{\textcolor[rgb]{0.81,0.36,0.00}{\textbf{#1}}}
\newcommand{\SpecialStringTok}[1]{\textcolor[rgb]{0.31,0.60,0.02}{#1}}
\newcommand{\StringTok}[1]{\textcolor[rgb]{0.31,0.60,0.02}{#1}}
\newcommand{\VariableTok}[1]{\textcolor[rgb]{0.00,0.00,0.00}{#1}}
\newcommand{\VerbatimStringTok}[1]{\textcolor[rgb]{0.31,0.60,0.02}{#1}}
\newcommand{\WarningTok}[1]{\textcolor[rgb]{0.56,0.35,0.01}{\textbf{\textit{#1}}}}
\usepackage{graphicx}
\makeatletter
\newsavebox\pandoc@box
\newcommand*\pandocbounded[1]{% scales image to fit in text height/width
  \sbox\pandoc@box{#1}%
  \Gscale@div\@tempa{\textheight}{\dimexpr\ht\pandoc@box+\dp\pandoc@box\relax}%
  \Gscale@div\@tempb{\linewidth}{\wd\pandoc@box}%
  \ifdim\@tempb\p@<\@tempa\p@\let\@tempa\@tempb\fi% select the smaller of both
  \ifdim\@tempa\p@<\p@\scalebox{\@tempa}{\usebox\pandoc@box}%
  \else\usebox{\pandoc@box}%
  \fi%
}
% Set default figure placement to htbp
\def\fps@figure{htbp}
\makeatother
\ifLuaTeX
\usepackage[bidi=basic,shorthands=off]{babel}
\else
\usepackage[bidi=default,shorthands=off]{babel}
\fi
\ifLuaTeX
  \usepackage{selnolig} % disable illegal ligatures
\fi
\setlength{\emergencystretch}{3em} % prevent overfull lines
\providecommand{\tightlist}{%
  \setlength{\itemsep}{0pt}\setlength{\parskip}{0pt}}
\usepackage{fontspec}
\defaultfontfeatures{Ligatures=TeX}
\usepackage{amsmath,amssymb}
\usepackage{booktabs}
\usepackage{array}
\usepackage{geometry}
\usepackage{enumitem}
\usepackage{xcolor}
\usepackage{tcolorbox}
\usepackage{fancyhdr}
\usepackage{titlesec}
\usepackage{multicol}

\geometry{margin=1.8cm, top=2.5cm, bottom=2cm}
\pagestyle{fancy}
\fancyhf{}
\rhead{\small Projektowanie badań — 2026/27}
\lhead{\small Zarządzanie informacją | ISI UJ}
\cfoot{\thepage}

\titleformat{\section}{\large\bfseries\color{blue!70!black}}{\thesection.}{0.5em}{}
\titleformat{\subsection}{\normalsize\bfseries}{\thesubsection.}{0.5em}{}
\titleformat{\subsubsection}{\small\bfseries\color{gray!70!black}}{\thesubsubsection.}{0.5em}{}
\titlespacing*{\section}{0pt}{1.5ex}{1ex}
\titlespacing*{\subsection}{0pt}{1ex}{0.5ex}

\setlist[itemize]{nosep, leftmargin=1.5em}
\setlist[enumerate]{nosep, leftmargin=1.5em}

\tcbset{boxrule=0.5pt, left=4pt, right=4pt, top=4pt, bottom=4pt, fonttitle=\bfseries\small}

\newtcolorbox{definicja}[1][]{colback=blue!5, colframe=blue!50!black, title={#1}}
\newtcolorbox{uwagabox}[1][]{colback=orange!10, colframe=orange!70!black, title={#1}}
\newtcolorbox{wzorbox}[1][]{colback=gray!5, colframe=gray!50!black, title={#1}}
\newtcolorbox{przykladbox}[1][]{colback=purple!5, colframe=purple!50!black, title={#1}}

\usepackage{bookmark}
\IfFileExists{xurl.sty}{\usepackage{xurl}}{} % add URL line breaks if available
\urlstyle{same}
% fallback for those not using the hyperref driver hyperxmp:
\makeatletter
\@ifundefined{xmpquote}{\newcommand{\xmpquote}[1]{#1}}{}
\makeatother
\hypersetup{
  pdftitle={Handout C04 --- Pozycje ankiety{,} indeks i odpowiedzi wielokrotne},
  pdfauthor={Marek Deja},
  pdflang={pl-PL},
  hidelinks,
  pdfcreator={LaTeX via pandoc}}

\title{Handout C04 --- Pozycje ankiety, indeks i odpowiedzi wielokrotne}
\usepackage{etoolbox}
\makeatletter
\providecommand{\subtitle}[1]{% add subtitle to \maketitle
  \apptocmd{\@title}{\par {\large #1 \par}}{}{}
}
\makeatother
\subtitle{Zarządzanie informacją --- część ilościowa --- 2026/27}
\author{Marek Deja}
\date{2026/27}

\begin{document}
\maketitle

\section{C04 --- 90 minut}\label{c04-90-minut}

10 min odtworzenie → 20 min konstrukt/skale → 25 min wspólne obliczenia
→ 25 min Z04 → 10 min oddanie. Pracuj na własnym ID; reguły są wspólne
dla S01--S20, a treść pytań pochodzi ze scenariusza.

\section{Od pojęcia do kolumny}\label{od-pojux119cia-do-kolumny}

Konstrukt → wskaźnik → pozycja z instrukcją → odpowiedź/zmienna →
uzasadniony indeks. Jedno pytanie mierzy jedno doświadczenie, ma
odbiorcę, okres odniesienia, kody i regułę braków. Pojedyncza pozycja
1--5 jest porządkowa; cyfry nie nadają równych odległości.

\section{Rekodacja i indeks}\label{rekodacja-i-indeks}

Pozycja 3 opisuje trudność. Odwracamy ją, aby wysoki wynik oznaczał
lepszą dostępność. Brak pozostaje brakiem, a kod 99 trzeba wcześniej
zamienić na NA.

\begin{Shaded}
\begin{Highlighting}[]
\NormalTok{badaniaZI}\SpecialCharTok{::}\FunctionTok{odwroc\_pozycje}\NormalTok{(}\FunctionTok{c}\NormalTok{(}\DecValTok{1}\NormalTok{, }\DecValTok{2}\NormalTok{, }\DecValTok{3}\NormalTok{, }\DecValTok{4}\NormalTok{, }\DecValTok{5}\NormalTok{, }\ConstantTok{NA}\NormalTok{))}
\end{Highlighting}
\end{Shaded}

\begin{verbatim}
## [1]  5  4  3  2  1 NA
\end{verbatim}

Średni indeks wymaga minimum 5/6 odpowiedzi. Zachowujemy kopię surowych
pozycji, aby nie odwrócić kolumny dwukrotnie. Odczytaj N ważnych
indeksów i liczbę pominiętych osób.

\begin{Shaded}
\begin{Highlighting}[]
\NormalTok{dane }\OtherTok{\textless{}{-}}\NormalTok{ badaniaZI}\SpecialCharTok{::}\FunctionTok{przygotuj\_ankiete}\NormalTok{(badaniaZI}\SpecialCharTok{::}\FunctionTok{dane\_przykladowe}\NormalTok{())}\SpecialCharTok{$}\NormalTok{dane}
\NormalTok{p }\OtherTok{\textless{}{-}}\NormalTok{ dane[, }\FunctionTok{paste0}\NormalTok{(}\StringTok{"pozycja\_"}\NormalTok{, }\DecValTok{1}\SpecialCharTok{:}\DecValTok{6}\NormalTok{)]}
\NormalTok{p}\SpecialCharTok{$}\NormalTok{pozycja\_3 }\OtherTok{\textless{}{-}}\NormalTok{ badaniaZI}\SpecialCharTok{::}\FunctionTok{odwroc\_pozycje}\NormalTok{(p}\SpecialCharTok{$}\NormalTok{pozycja\_3)}
\NormalTok{i }\OtherTok{\textless{}{-}}\NormalTok{ badaniaZI}\SpecialCharTok{::}\FunctionTok{indeks\_ankiety}\NormalTok{(p, }\AttributeTok{minimum =} \DecValTok{5}\NormalTok{)}
\FunctionTok{c}\NormalTok{(}\AttributeTok{wazne =} \FunctionTok{sum}\NormalTok{(}\SpecialCharTok{!}\FunctionTok{is.na}\NormalTok{(i)), }\AttributeTok{braki =} \FunctionTok{sum}\NormalTok{(}\FunctionTok{is.na}\NormalTok{(i)))}
\end{Highlighting}
\end{Shaded}

\begin{verbatim}
## wazne braki
##   147     1
\end{verbatim}

Indeks jest przybliżonym pomiarem złożonym. Samo uśrednienie ani alfa
Cronbacha nie dowodzą trafności. Osoby z pięcioma odpowiedziami mogą
mieć wynik oparty na innym zestawie pozycji.

\section{Kanały --- respondent jes
mianownikiem}\label{kanaux142y-respondent-jest-mianownikiem}

0 = niezaznaczenie, 1 = zaznaczenie, NA = brak. Reguła kursu obejmuje
tylko kompletne cztery pola. Jedna osoba może wybrać wiele opcji, więc
suma procentów może przekraczać 100\%.

\begin{Shaded}
\begin{Highlighting}[]
\NormalTok{badaniaZI}\SpecialCharTok{::}\FunctionTok{odpowiedzi\_wielokrotne}\NormalTok{(dane[, }\FunctionTok{paste0}\NormalTok{(}\StringTok{"kanal\_"}\NormalTok{, }\DecValTok{1}\SpecialCharTok{:}\DecValTok{4}\NormalTok{)])}
\end{Highlighting}
\end{Shaded}

\begin{verbatim}
##     opcja liczba mianownik procent_respondentow pominiete_niepelne
## 1 kanal_1     91       144             63.19444                  4
## 2 kanal_2     76       144             52.77778                  4
## 3 kanal_3     59       144             40.97222                  4
## 4 kanal_4     79       144             54.86111                  4
\end{verbatim}

\section{Z04 --- dowody i punktacja}\label{z04-dowody-i-punktacja}

Własny słownik/pytanie/skale/kody --- 3 pkt. Pozycja 3 i minimum 5/6 ---
3 pkt. Kanały/liczebności/procent respondentów/mianownik --- 2 pkt.
Wyjaśnienie pozycji i indeksu z ograniczeniem pomiaru --- 2 pkt.

Pliki: \texttt{zadania/z04/analiza.R}, \texttt{odpowiedzi.md},
\texttt{wyniki/slownik\_pomiaru.csv}, \texttt{indeks.csv},
\texttt{kanaly.csv}. Zapisz linki w MD, uruchom skrypt od początku i
sprawdź odbiór przez funkcje R. Pełne poziomy: \texttt{rubryka("Z04")}.

\section{Notatki}\label{notatki}

Co mierzę? \ldots\ldots\ldots\ldots\ldots\ldots\ldots\ldots\ldots{}
Dlaczego pozycja 3 ma inny kierunek?
\ldots\ldots\ldots\ldots\ldots\ldots\ldots.

N ważnych indeksów: \ldots\ldots\ldots. Mianownik kanałów:
\ldots\ldots\ldots. Niepełne odpowiedzi: \ldots\ldots\ldots.

Jedno ograniczenie pomiaru:
\ldots\ldots\ldots\ldots\ldots\ldots\ldots\ldots\ldots\ldots\ldots\ldots\ldots\ldots\ldots\ldots\ldots\ldots\ldots\ldots\ldots\ldots..

\end{document}
